\chapter{Bayes Nets as Models of Rationality}
\label{sec:p16-rationality-4-2}

Hartmann's chapter distinguishes the general Bayesian framework from
the particular models built within it. A probability calculus does not
select the relevant propositions, their dependencies, the source of a
report, or the task an agent is trying to solve. Bayes nets make those
modeling choices explicit and permit their consequences to be calculated
\cite[chapter~4.2, pp.~253--264]{Hartmann2021BayesNets}.
The chapter's examples concern confirmation, testimony and argumentation.
Their arrows primarily represent a joint probability law; their causal
vocabulary is explicitly heuristic, not an intervention semantics
\cite[note~4, pp.~262--263]{Hartmann2021BayesNets}.

This treatment follows that account through the three elementary networks,
partially reliable reports, the variety-of-evidence comparison, the
no-alternatives and no-reason-for arguments, and learning conditionals.
It keeps graphical representation, normative assessment and psychological
interpretation separate. The numerical illustrations below are stipulated
models, not measurements of actual witnesses or scientific communities.
Two selected primary empirical records are identified separately.

\section{What a graph represents and what its probabilities add}
\label{sec:p16-rat42-graphical-representation}

\begin{hbconvention}{def:p16-rat42-notation}{Variables, values and the common probability law}
Following the chapter's note~2, italic letters such as \(A,B,H\)
denote variables; roman \(\mathrm A,\neg\mathrm A\) and
\(\mathrm H,\neg\mathrm H\) denote their binary values or the
corresponding events. Thus \(P(\mathrm A,\mathrm B)\) abbreviates
\(P(\mathrm A\land\mathrm B)\), not a new probability operation.
The multivalued variable \(Y\) has value events \(\mathrm Y_i\).
The graphical tables below have finite value sets; the no-reason-for
argument explicitly also permits a countable number-of-arguments variable.
Each calculation uses one specified probability law \(P\).
Conditional probabilities are asserted only on
events of positive probability. We retain the source's shorthand
\(\bar x=1-x\) for a probability parameter, distinct from event
negation. A posterior law is denoted \(P'\).
\end{hbconvention}

\begin{hbdefinition}{def:p16-rat42-bayes-net}{A finite Bayes net and the parental Markov condition}
A Bayes net has a finite DAG whose nodes are variables
\(A_1,\ldots,A_n\), together with a normalized nonnegative local
probability table \(k_i(a_i\mid a_{pa(A_i)})\) at each node,
specified for every parent assignment. For a root the parent set
\(pa(A_i)\) is empty, so its table is a prior mass function.
A root has no incoming arrow; an isolated node is also a root.
A descendant is reached by a directed path of positive length.

The \emph{parental Markov condition} (PMC) states that each node's
variable is jointly independent of its nondescendant nonparent
variables, conditional on its parents. The node itself is excluded
from its nondescendants. This is joint conditional independence,
not merely a list of pairwise independence claims. For finite
variables, \(A\perp B\mid C\) means
\[
 P(a,b\mid c)=P(a\mid c)P(b\mid c)
 \quad\text{whenever }P(c)>0.
\]
It is equivalent to \(P(a\mid b,c)=P(a\mid c)\) on the
positive \((b,c)\) strata. An arrow alone specifies neither a
numerical probability nor a counterfactual or intervention effect.
\end{hbdefinition}

\begin{hbtheorem}{thm:p16-rat42-product}{proved=annex:p16.rat42-product}{Product representation and its exact parameter count}
The local tables define a normalized joint law
\[
 P(a_1,\ldots,a_n)=\prod_{i=1}^n k_i(a_i\mid a_{pa(A_i)}).
\]
It satisfies PMC. Conversely a law satisfying PMC factors in this
way, with its ordinary conditionals on positive parent strata;
the unspecified tables on null strata can be filled arbitrarily.
For binary variables the unrestricted joint has \(2^n-1\) free
parameters, while unrestricted local tables on the given DAG have
\(\sum_i2^{|pa(A_i)|}\). Sparse graphs can reduce this count
substantially; a fully connected topologically ordered DAG does not.
Every finite marginal is obtained by summing this joint law, and
every positive-event conditional by summing and normalizing it.
\end{hbtheorem}

The source's virtues are representational economy, systematic probability
calculation, and an accessible graphical theory. Economy depends on which
independences are defensible. Efficient implementations exploit the graph,
but a compact specification is not by itself a guarantee that every
inference query is inexpensive. No particular cognitive algorithm follows
from the product representation.

\begin{hblemma}{thm:p16-rat42-semigraphoid}{proved=annex:p16.rat42-semigraphoid}{The conditional-independence operations used by the chapter}
For finite variable tuples, conditional independence satisfies symmetry,
decomposition, weak union and contraction:
\[
\begin{aligned}
 A\perp B\mid C&\Longrightarrow B\perp A\mid C,\\
 A\perp(B,D)\mid C&\Longrightarrow A\perp B\mid C,\\
 A\perp(B,D)\mid C&\Longrightarrow A\perp B\mid(C,D),\\
 A\perp B\mid C,\quad A\perp D\mid(B,C)
     &\Longrightarrow A\perp(B,D)\mid C.
\end{aligned}
\]
No assertion of an intersection axiom on arbitrary zero-containing
distributions is included.
\end{hblemma}

The chapter illustrates these relations with its figures~4.2.1--4.2.3:
\[
 \begin{array}{c|c}
  \text{source graph}&\text{graph-guaranteed independence}\\\hline
  A\leftarrow C\to B&A\perp B\mid C\\
  B\to C\to A&A\perp B\mid C\\
  A\to C\leftarrow B&A\perp B
 \end{array}
\]
The first two describe the same independence constraint despite different
orientations. Their equality does not identify a causal direction.
The collider is different: conditioning on its common effect can make
the parent variables dependent, rather than screening them off.

\begin{hbtheorem}{thm:p16-rat42-three-networks}{proved=annex:p16.rat42-three-networks}{Equivalent noncolliders and a calculated collider contrast}
A finite law with \(A\perp B\mid C\) has both the common-parent
and chain factorizations displayed above. In the binary collider let
\[
\begin{gathered}
 P(\mathrm A)=a,\quad P(\mathrm B)=b,\quad A\perp B,\\
 P(\mathrm C\mid\mathrm A,\mathrm B)=\alpha,\quad
 P(\mathrm C\mid\mathrm A,\neg\mathrm B)=\beta,\\
 P(\mathrm C\mid\neg\mathrm A,\mathrm B)=\gamma,\quad
 P(\mathrm C\mid\neg\mathrm A,\neg\mathrm B)=\delta.
\end{gathered}
\]
Then \(P(\mathrm A,\mathrm B,\mathrm C)=ab\alpha\) and
\(P(\mathrm A,\neg\mathrm B,\mathrm C)=a\bar b\beta\).
For \(a=b=1/2\) and
\((\alpha,\beta,\gamma,\delta)=(9/10,4/5,7/10,1/10)\),
\[
 P(\mathrm A\mid\mathrm B,\mathrm C)=9/16,
 \qquad P(\mathrm A\mid\mathrm C)=17/25.
\]
Thus this collider's parents become dependent given \(\mathrm C\).
If instead all four conditional parameters are equal to \(2/5\),
the parents remain conditionally independent. Collider structure
permits the dependence; it does not force it for every table.
\end{hbtheorem}

For larger graphs the source points to d-separation rather than enumerating
all independence consequences manually. The exact active-path criterion
and its complete finite soundness proof are
\cref{def:p16-active-path,thm:p16-finite-global-markov}; they apply to
the normalized finite product law above without requiring causal equations.
They supply independence implications guaranteed by factorization.
Additional accidental independences can hold for a particular numerical
law, as the constant collider table demonstrates. Inferring the entire
graph from observed independences needs further assumptions, not just PMC.

\section{The Bayesian framework is not the particular model}
\label{sec:p16-rat42-framework-and-model}

Hartmann separates a static requirement, representing coherent degrees of
belief by a probability law over a selected Boolean algebra, from a dynamic
requirement, revising that law when information arrives. For example,
credence \(.3\) in tomorrow's rain requires \(.7\) in its negation
if they are genuinely complementary events under the same law. Selecting
which propositions to include precedes that coherence calculation.
\cite[Sec.~3.1, pp.~255--256]{Hartmann2021BayesNets}

For certain evidence \(\mathrm E\) with positive prior probability,
conditionalization sets \(P'(X)=P(X\mid\mathrm E)\) for every event
\(X\). The chapter cites pragmatic Dutch-book arguments and epistemic
accuracy arguments as distinct normative foundations. Its account of
conservativity supplies another proposed principle: satisfy the new
information while changing the old beliefs as little as possible.
Specifying what counts as little change is an additional choice.
The local finite result below makes one such choice precise; it is not
evidence that people actually optimize that quantity.

\begin{hbdefinition}{def:p16-rat42-kl}{Directed divergence and a learning constraint}
For mass functions \(Q,P\) on the same finite outcome space, set
\[
 D(Q\|P)=\sum_{\omega:Q(\omega)>0}
                  Q(\omega)\log\frac{Q(\omega)}{P(\omega)}.
\]
The value is \(+\infty\) if a positive \(Q\)-mass has zero
\(P\)-mass; zero \(Q\)-terms contribute zero. This is the
Kullback--Leibler divergence, directed from proposed posterior to prior,
not a symmetric distance or a metric. A learning constraint is a
specified condition on \(Q\), such as \(Q(\mathrm E)=1\) or
\(Q(\mathrm C\mid\mathrm A)=t\). The latter also requires
\(Q(\mathrm A)>0\).
\end{hbdefinition}

\begin{hbtheorem}{thm:p16-rat42-kl-update}{proved=annex:p16.rat42-kl-update}{The finite conservativity calculation and its direction}
If \(P(\mathrm E)>0\), the unique minimizer of \(D(Q\|P)\)
subject to \(Q(\mathrm E)=1\) is \(Q=P(\cdot\mid\mathrm E)\).
For a strictly positive law on the four \((A,C)\) assignments, put
\(s=P(\mathrm A)\), \(u=P(\mathrm C\mid\mathrm A)\), and
\(v=P(\mathrm C\mid\neg\mathrm A)\). For the constraint
\(Q(\mathrm C\mid\mathrm A)=t\in[0,1]\), define
\[
 d=t\log(t/u)+(1-t)\log((1-t)/(1-u)),
 \qquad
 s'=\frac{s e^{-d}}{\bar s+s e^{-d}},
\]
using zero-term conventions. The unique minimizer has
\(Q(\mathrm A)=s'\), \(Q(\mathrm C\mid\mathrm A)=t\),
and \(Q(\mathrm C\mid\neg\mathrm A)=v\).
When \(t=1\), this law is also
\(P(\cdot\mid\neg\mathrm A\lor\mathrm C)\).
The mathematical coincidence for this hard constraint does not establish
that all indicative conditionals mean material implications.
\end{hbtheorem}

The general framework does not choose a graph, a reliability parameter,
a likelihood, or this particular divergence. Hartmann compares it with
Newtonian mechanics: a framework and general laws do not specify the
system to be modeled or its initial conditions. Particular models apply
the framework, expose assumptions, organize bookkeeping, solve questions,
and yield predictions that can be compared with evidence. Model simplicity
and fidelity to details can conflict. Combining variables into a single
macro-variable or deleting an alleged weak dependence requires a substantive
justification; it is not licensed merely by a cleaner diagram.
\cite[Secs.~3.1--3.2, pp.~256--257]{Hartmann2021BayesNets}

\begin{hbmethod}{def:p16-rat42-modeling-method}{From a reasoning question to an assessable Bayes-net model}
\begin{enumerate}
\item State the question, the agent's information, and the target:
      a posterior, confirmation comparison, argument assessment, or
      descriptive prediction. Distinguish a normative target from an
      observed response measure.
\item Specify variable meanings and mutually exclusive, exhaustive values.
      Include uncertainty about a source when it matters; do not silently
      equate a report with the reported fact.
\item State the graph's joint independence assumptions and give the
      particular scenario's reasons for them. Identify alternative
      dependencies, omitted variables and plausible coarser models.
\item Specify every local table and its provenance: elicited credences,
      externally estimated rates, fixed task information, or fitted
      parameters. Record zeros and the learning constraints.
\item Form the joint product, sum unwanted variables, and normalize only
      positive evidence events. Inspect limiting and contrary cases.
\item Evaluate the resulting claim at its proper level. An internally
      coherent posterior is conditional on the model. A descriptive
      prediction needs task-matched evidence and comparison with alternatives.
\end{enumerate}
This reconstructs the chapter's framework/model distinction and its
functions of models; it does not automate variable selection or supply
a universal graph-learning rule.
\end{hbmethod}

\section{Confirmation and partially reliable reports}
\label{sec:p16-rat42-reliability-confirmation}

\begin{hbdefinition}{def:p16-rat42-confirmation}{Confirmation relative to an agent's law}
For \(0<P(\mathrm H)<1\) and \(P(\mathrm E)>0\),
\(\mathrm E\) confirms \(\mathrm H\) if
\(P(\mathrm H\mid\mathrm E)>P(\mathrm H)\), disconfirms it if
the inequality is reversed, and is irrelevant if equality holds.
This is a comparison of credences, not deductive proof or refutation.
Let \(h=P(\mathrm H)\), \(p=P(\mathrm E\mid\mathrm H)\),
and \(q=P(\mathrm E\mid\neg\mathrm H)\). The source uses
the inverse likelihood ratio \(x=q/p\) when \(p>0\).
\end{hbdefinition}

\begin{hblemma}{thm:p16-rat42-confirmation}{proved=annex:p16.rat42-confirmation}{The posterior and the likelihood comparison}
Under these assumptions,
\[
 P(\mathrm H\mid\mathrm E)=\frac{hp}{hp+\bar hq}
     =\frac{h}{h+\bar h x}\quad(p>0).
\]
Confirmation, disconfirmation and irrelevance are respectively
\(p>q\), \(p<q\), and \(p=q\). If \(p=0<q\), the
posterior is zero but \(q/p\) is undefined; if \(q=0<p\),
it is one. If \(p=q=0\), the evidence cannot be conditioned on.
\end{hblemma}

Hartmann starts with error rates: false positives are
\(fp=P(\mathrm E\mid\neg\mathrm H)=q\), and false negatives
are \(fn=P(\neg\mathrm E\mid\mathrm H)=1-p\).
Those rates can sometimes be estimated for an instrument, but are often
unknown for testimony. A larger graph can unpack them into interpretable
uncertainties rather than treating them as unexplained numbers.
\cite[Sec.~3.3, pp.~257--259]{Hartmann2021BayesNets}

\begin{hbdefinition}{def:p16-rat42-source-model}{The chapter's reliable-or-randomizing witness}
Let \(H\) be the hypothesis that Paul is the murderer, \(Rel\)
the witness's reliability, and \(Rep\) the witness's report.
The graph is \(H\to Rep\leftarrow Rel\), with independent roots
\(P(\mathrm H)=h\), \(P(\mathrm{Rel})=r\). A reliable source
always reports the truth; an unreliable one reports positively with
probability \(a\), independently of the truth. Thus
\[
\begin{array}{c|cc}
 &\mathrm{Rel}&\neg\mathrm{Rel}\\\hline
 P(\mathrm{Rep}\mid\mathrm H,\cdot)&1&a\\
 P(\mathrm{Rep}\mid\neg\mathrm H,\cdot)&0&a
\end{array}
\]
The parameters are the modeled agent's assumptions, not measured
properties of an actual witness. Dishonesty, honest error, copying,
and changing reliability are not all represented by this binary split.
\end{hbdefinition}

\begin{hbtheorem}{thm:p16-rat42-witness}{proved=annex:p16.rat42-witness}{One report, two sources, and the failure of a universal variety claim}
For the one-report model,
\[
 p=r+a\bar r,\quad q=a\bar r,\qquad
 P(\mathrm H\mid\mathrm{Rep})=
                   \frac{h(r+a\bar r)}{hr+a\bar r}
\]
whenever the denominator is positive. For \(0<h<1\) and
\(r>0\), a positive report confirms \(\mathrm H\).
For \(r=0\), it is irrelevant if \(a>0\); for \(r=1\),
it makes the hypothesis certain if the report has positive probability.

Compare two positive reports, with the same \(h,r,a\).
In the separate-source model the roots \(H,Rel_1,Rel_2\) are
mutually independent; \(Rep_j\) has parents \(H,Rel_j\) and
uses the one-report table. The reports are independent given
\(H,Rel_1,Rel_2\). In the shared-source model the independent
roots \(H,Rel\) parent both reports, each using that same table;
the reports are independent given \(H,Rel\). The respective likelihoods are
\[
 (p_I,q_I)=((r+a\bar r)^2,(a\bar r)^2),
 \qquad (p_D,q_D)=(r+a^2\bar r,a^2\bar r).
\]
For \(0<h,r,a<1\), the shared-reliability model confirms more
strongly precisely when
\[
        a<\tfrac12,
        \qquad r<\frac{1-2a}{2(1-a)}.
\]
At \(h=1/2,r=a=1/10\), the posteriors are
\(361/442\) for separate reliability and \(109/118\) for shared
reliability. At \(a=1/2\) and \(0<r<1\), the separate-source
model confirms more strongly. Thus dependence can help or hurt;
equal one-report marginals do not settle the two-report comparison.
\end{hbtheorem}

The source's figures~4.2.5--4.2.6 motivate this contrast and refer to
Bovens and Hartmann for its analysis. The displayed derivation is supplied
locally. The separate reports are not unconditionally independent: both
carry information about the shared \(H\). The common-reliability model
is also not a model of two literal copies of one report. Its conditionally
independent randomizations are part of the comparison. This explains why
the result is compatible with the duplicated-report calculation in
\cref{thm:p15-bayes-flag-calculation}.

The chapter connects these models to the variety-of-evidence thesis,
Landes's examination of its conditions, and applications involving learning
source reliability. It also sketches analogue confirmation: observing
a testable phenomenon predicted by \(T'\) may bear on an untestable
analogous prediction of \(T\) only through additional common structural
assumptions. An analogy alone supplies no likelihood for the target.
The legal applications likewise concern intertwined testimony and source
credibility. These are the chapter's attributed application directions,
not a claim that a single reliability graph proves every such inference.
\cite[Sec.~3.3, pp.~258--259]{Hartmann2021BayesNets}

\section{Argumentation as learning: validity and probabilistic force}
\label{sec:p16-rat42-argument-learning}

In the chapter's reconstruction, an agent has beliefs about premises and
a conclusion. Another person attempts to alter the agent's beliefs about
the premises; the agent then revises the joint law. An argument has
probabilistic force for this agent when the revised probability of its
conclusion exceeds the old probability. This is a relative confirmation
criterion, not a claim that every persuasive argument is sound or that
truth is whatever an audience finds convincing. Deductive relations
constrain the allowed joint distributions, while the content and the
learning procedure determine the actual change.
\cite[Sec.~4, pp.~259--260]{Hartmann2021BayesNets}

\begin{hbtheorem}{thm:p16-rat42-validity}{proved=annex:p16.rat42-validity}{Certain material premises and specified partial learning}
If the event \(\mathrm A\land(\neg\mathrm A\lor\mathrm C)\)
has positive probability, learning both material modus-ponens premises
by conditionalization yields
\[
 P'(\mathrm C)=P(\mathrm C\mid\mathrm A,
                                  \neg\mathrm A\lor\mathrm C)=1.
\]
For partial learning, suppose the material major premise already has
probability one, \(0<s=P(\mathrm A)<1\), and
\(t=P(\mathrm C\mid\neg\mathrm A)\). Change \(P(\mathrm A)\)
to \(s'>s\) while preserving the conditionals of \(C\) on
\(A\) and \(\neg A\), as in a binary Jeffrey update. Then
\[
 P'(\mathrm C)-P(\mathrm C)=(s'-s)(1-t)\geq0.
\]
The increase is strict only if \(t<1\); if \(s'<1\) and
\(t<1\), the posterior remains below one.

Denying the antecedent is invalid deductively but can have either
probabilistic direction. With \(P(\mathrm A)=1/2\), choosing
\((P(\mathrm C\mid\mathrm A),P(\mathrm C\mid\neg\mathrm A))
=(.9,.2)\) makes learning \(\neg\mathrm A\) and the material
conditional raise \(P(\neg\mathrm C)\) from \(.45\) to \(.8\).
Choosing instead \((.1,.8)\) lowers it from \(.55\) to \(.2\).
\end{hbtheorem}

This makes the chapter's valid/invalid contrast precise without assigning
an automatic increase to every unspecified change in uncertain premises.
Validity guarantees a true conclusion when the premises are jointly
certain under fixed meanings. With partial premise information, which
conditionals are retained and which probabilities change must be stated.
The fixed material validity result in
\cref{thm:p19-conditional-benchmark} does not contradict observed suppression
or a changed interpretation of an indicative conditional.

\section{No alternatives and no reasons: the hidden common variable}
\label{sec:p16-rat42-search-arguments}

\begin{hbdefinition}{def:p16-rat42-common-parent}{The number-of-candidates model}
Let \(Y\) take mutually exclusive values \(\mathrm Y_i\), with
probabilities \(y_i\) summing to one. Let
\(h_i=P(\mathrm H\mid\mathrm Y_i)\) and
\(f_i=P(\mathrm F\mid\mathrm Y_i)\) on positive \(y_i\);
on zero strata these symbols name stipulated local table entries,
not conditional probabilities on null events. The graph
\(H\leftarrow Y\to F\) stipulates \(H\perp F\mid Y\),
equivalently the joint table
\[
       P(\mathrm H,\mathrm F,\mathrm Y_i)=y_i h_i f_i,
\]
with the other three \((H,F)\) cells obtained by complementing
\(h_i\) and \(f_i\). The index set may be finite or countable;
all resulting series converge absolutely. Assume
\(P(\mathrm F)=\sum_i y_i f_i>0\) when conditioning on \(F\).
\end{hbdefinition}

\begin{hblemma}{thm:p16-rat42-covariance}{proved=annex:p16.rat42-covariance}{The exact confirmation identity for the common-parent model}
For this model,
\[
 P(\mathrm H\mid\mathrm F)-P(\mathrm H)
 =\frac{\displaystyle
   \frac12\sum_{i,j}y_i y_j(h_i-h_j)(f_i-f_j)}
        {\displaystyle\sum_i y_i f_i}.
\]
The formula applies also to countably many values of \(Y\).
The conditional-independence assumption is what makes the mixed
expectation equal to \(\sum_i y_i h_i f_i\).
\end{hblemma}

\subsection{The no-alternatives argument}
\label{sec:p16-rat42-no-alternatives}

Hartmann's scientist has a theory \(H\) satisfying desirable constraints
but cannot yet test its predictions. Despite substantial search, the
community has found no alternative satisfying those same constraints.
The proposed conclusion is a reason in favor of \(H\), not its deductive
truth. The chapter uses the discussion of string theory as its historical
motivation; this is a reconstruction of that 2021 argument, not a current
forecast about empirical physics.
\cite[Sec.~4.1, pp.~260--261]{Hartmann2021BayesNets}

In this application \(\mathrm F\) means no suitable alternative has yet
been found, and \(\mathrm Y_i\) means exactly \(i\) suitable alternatives
exist. The source uses a finite maximum \(N\). Its assumptions are:
\[
\begin{array}{ll}
 \mathrm{NAA1}:&0\leq y_i<1,\quad \sum_{i=0}^N y_i=1;\\
 \mathrm{NAA2}:&h_i\text{ is nonincreasing in }i;\\
 \mathrm{NAA3}:&f_i\text{ is nonincreasing in }i;\\
 \mathrm{NAA4}:&H\perp F\mid Y.
\end{array}
\]
The restriction \(y_i<1\) keeps uncertainty about the count;
it is not a consequence of coherence alone. More alternatives can
lower the hypothesis's credence only under the specified comparison of
suitable candidates. More alternatives make search failure less likely
only if search is sensitive to them. The screening-off assumption says
the count carries all relevant dependence between the theory and the
search outcome. It is a modeling premise, not something guaranteed by
the variable's name. A lack of logical entailment between \(F\) and
\(H\) does not rule out direct probabilistic dependence between them.

\begin{hbtheorem}{thm:p16-rat42-naa}{proved=annex:p16.rat42-naa}{The no-alternatives confirmation criterion}
Under NAA1--NAA4 and \(P(\mathrm F)>0\), \(F\) confirms \(H\)
exactly when there are \(i>j\) with
\[
           y_i y_j>0,\qquad f_i<f_j,
                          \qquad h_i<h_j.
\]
The posterior change is otherwise zero, not negative, under these
same monotonicity assumptions.
For \(y_0=y_1=1/2\),
\((h_0,h_1)=(4/5,1/5)\), and
\((f_0,f_1)=(4/5,1/5)\), the prior is \(1/2\),
the probability of \(F\) is \(1/2\), and the posterior is \(17/25\).
\end{hbtheorem}

The theorem exposes the argument's vulnerable premises. If the agent is
certain of the number of alternatives, no count uncertainty is available
to mediate confirmation. The philosophical objection that appropriate
alternatives may be unlimited therefore challenges the finite count model
and its priors, rather than being settled by this theorem. A restricted
class of suitable theories must be justified, not silently treated as
all logically possible competitors.

Search difficulty is a further competing explanation. Adding a variable
\(D\) for that difficulty need not remove all confirmation, but it can
make the marginal \(f_i\)'s constant, leaving \(F\) uninformative about
\(Y\) and hence \(H\). For example, a search procedure that never
reports a found alternative has \(f_i=1\) for every \(i\); then
\(P(\mathrm H\mid\mathrm F)=P(\mathrm H)\) whatever the count prior.
Less extreme difficulty can merely weaken, rather than eliminate, the
signal. Finally, community agreement that no alternative has been found
is a different proposition from no alternative existing. An adequate
search process, agreed suitability standards and reliable reporting are
needed to justify the \(f_i\)'s.

\subsection{The no-reason-for argument}
\label{sec:p16-rat42-no-reason-for}

The source's contrasting example examines someone who searches for good
arguments for God's existence and finds none. It asks whether that search
outcome can constitute a reason against the hypothesis. The reconstructed
\(H\) can be any hypothesis: \(\mathrm F\) means no good argument for
it has yet been found, and \(\mathrm Y_i\) means exactly \(i\) good
arguments for it exist. Here the count can range over all nonnegative
integers. The model assumptions are:
\[
\begin{array}{ll}
 \mathrm{NRF1}:&0\leq y_i<1,\quad\sum_i y_i=1;\\
 \mathrm{NRF2}:&h_i\text{ is nondecreasing in }i;\\
 \mathrm{NRF3}:&f_i\text{ is nonincreasing in }i;\\
 \mathrm{NRF4}:&H\perp F\mid Y.
\end{array}
\]
Unlike NAA2, NRF2 makes more supporting arguments raise the hypothesis's
credence. This is a substantive relation between good reasons and truth,
not a count of rhetorically persuasive but unsound arguments. NRF3 models
a search more likely to succeed when more good arguments are available.

\begin{hbtheorem}{thm:p16-rat42-nrf}{proved=annex:p16.rat42-nrf}{The no-reason-for disconfirmation criterion}
Under NRF1--NRF4 and \(P(\mathrm F)>0\), \(F\) disconfirms \(H\)
exactly when there are \(i>j\) with
\[
            y_i y_j>0,\qquad f_i<f_j,
                           \qquad h_i>h_j.
\]
The posterior change is otherwise zero, not positive, under these
same assumptions. For \(y_0=y_1=1/2\),
\((h_0,h_1)=(1/5,4/5)\), and
\((f_0,f_1)=(4/5,1/5)\), the prior is \(1/2\) and the
posterior is \(8/25\).
\end{hbtheorem}

This is a conditional vindication of the argument under an explicit
search-and-reasons model. It does not make every failure to discover an
argument evidence against a hypothesis. An uninformative search, fixed
count, or missing relationship between argument availability and truth
can remove the conclusion. The model does not settle the theological
example's premises by assigning them labels.

\section{Indicative conditionals and the interpretation of new information}
\label{sec:p16-rat42-indicative-learning}

The source's section~4.3 returns to a limit of the material-premise
illustration. An indicative sentence of the form ``If A, then C'' need
not assert the truth-functional proposition \(\neg\mathrm A\lor\mathrm C\).
Treating it that way makes conditioning easy but does not settle its
meaning. Some accounts do not represent its content by a single event
at all. Hartmann therefore sketches a more general procedure: translate
what was learned into a constraint such as
\(P'(\mathrm C\mid\mathrm A)=1\), then choose a conservative posterior
subject to that constraint. His cited Eva--Hartmann--Rafiee Rad proposal
distinguishes interpretation of evidence from the update implementing it.
\cite[Sec.~4.3, p.~262]{Hartmann2021BayesNets}

The directed-KL calculation above shows both what this proposal can do
and what it leaves to interpretation. It handles constraints on
conditional probabilities, not only certainty in an event. A softer
constraint \(t<1\) is not ordinary conditioning on the material
conditional; it can allow positive mass for
\(\mathrm A\land\neg\mathrm C\). Different distances, additional
constraints, or a different account of testimony may produce different
posteriors. A perfectly reliable informant is an explicit idealization,
not a default description of every speaker.

The chapter also points to Bayes-net models of disablers. A disabler adds
a relevant condition on which the original dependence or likelihood can
change. For instance, a report that a switch normally turns a lamp on
must be reconsidered when a power-failure variable is learned. Introducing
that variable changes the modeled information and possibly the original
conditional's intended background. It does not create a counterexample
to fixed material modus ponens. The empirical suppression comparison in
\cref{sec:p19-conditional-reasoning-suppression} is therefore an adjacent
task with its own premise interpretations, not proof that every disabler
instantiates this particular distance-minimizing account.

\section{Cognitive interpretation: two bounded empirical comparisons}
\label{sec:p16-rat42-cognitive-evidence}

The psychological question is not whether the product rule is true but
which representation, prior, information and response objective a person
uses. A normative network can be tested as a computational-level account
of judgments without claiming a particular mental algorithm. Aggregate
agreement, individual prediction, mechanism identification and population
generalization are different evidential goals, as
\cref{def:p19-benchmark,def:p19-evidence-record} make explicit.
The following two primary studies are the finite empirical comparison
set for this treatment; their studies are not pooled into a universal
estimate of Bayesian reasoning.

\subsection{Expertise and trustworthiness in an appeal to an expert}
\label{sec:p16-rat42-expert-study}

Harris, Hahn, Madsen and Hsu distinguish expertise from trustworthiness:
an expert knows the fact, while a trustworthy person reports what they
believe. Their simplified model assumes an untrustworthy expert reports
the opposite, and an uninformed source randomizes fairly. This is not
Hartmann's reliable-or-randomizing binary source: informed systematic
deceit has been separated from ignorance.
\cite[accepted manuscript, pp.~12--16]{HarrisHahnMadsenHsu2016Expert}

\begin{hblemma}{thm:p16-rat42-expertise-trust}{proved=annex:p16.rat42-expertise-trust}{The distinct roles of expertise and honesty}
Let independent roots \(H,Exp,T\) have probabilities
\(h,e,t\). Given \(Exp\), an honest source reports the truth and
a dishonest source its opposite; given \(\neg Exp\), either
source reports positively with probability \(1/2\). Then
\[
 p=et+(1-e)/2,\quad q=e(1-t)+(1-e)/2,
 \quad p-q=e(2t-1).
\]
For \(0<h<1\), a positive report confirms, is irrelevant, or
disconfirms the hypothesis according to the sign of \(e(2t-1)\).
For \(e=4/5\), \(t=9/10\) gives likelihoods \((41/50,9/50)\),
whereas \(t=1/10\) reverses them. More informed systematic deception
need not be more favorable evidence for the stated hypothesis.
\end{hblemma}

\paragraph{Task, sample and design.}
Experiment~1 analyzed 84 unpaid online volunteers, 29 male and 55 female,
aged 18--80, median 20, recruited through an online experiment listing.
Each judged four dialogues concerning fictional medical claims.
Expertise was doctor versus musician; trustworthiness was friend versus
enemy of the argument proponent. Both were within-person factors.
Other experts' opinions were a three-level between-person factor:
no mention, agreement or disagreement. A Latin-square assignment varied
which content instantiated each within-person combination.
Participants used unnumbered sliders, recorded on a 0--100 scale,
to judge what the dialogue's recipient \emph{should} believe and to
estimate expertise and trust before and after the report. These are
elicited judgments about a recipient, not observed beliefs of that
fictional person or medical measurements.
\cite[accepted manuscript, pp.~16--20]{HarrisHahnMadsenHsu2016Expert}

\paragraph{Comparison, findings and uncertainty.}
The local conditional table was fixed by the researchers. Expertise and
trust priors were elicited; the hypothesis prior was fitted separately
for the three other-expert conditions, giving three fitted parameters.
The model accounted for a reported 89\% of variance across condition-mean
convincingness ratings. Across the three response types within individuals,
the reported variance account was 38\%, not 89\%. The group statistic
must not be presented as individual predictive accuracy. The authors
reported a systematic model/data interaction involving expertise and
others' opinions, \(F(2,70)=3.81,p=.027\), after excluding predictions
of certainty. The mismatch was particularly visible for a trustworthy
doctor disagreeing with other experts. Their figures use 95\% confidence
intervals; no numerical interval for the reported variance accounts is
supplied here or invented.
\cite[accepted manuscript, pp.~20--25]{HarrisHahnMadsenHsu2016Expert}

Experiment~2 retained 57 unpaid online volunteers, 23 male and 34 female,
aged 16--54, median 23, after excluding two respondents under 16 and one
reporting age 99. It reworded expertise as whether the source knows the
answer and elicited the hypothesis prior rather than fitting it.
The three condition sizes were 18, 19 and 20. With no fitted prior
parameters, but the same stipulated conditional table, it reported 94\%
variance account for group convincingness and again 38\% at the individual
level across response types. Among 50 participants for whom the
four-convincingness-response correlation comparison was possible, the
full model fit best for 27, a model omitting trust for 15, and a model
omitting expertise for 8. Four points per individual provide only weak
discrimination among process accounts.
\cite[accepted manuscript, pp.~25--31]{HarrisHahnMadsenHsu2016Expert}

The authors explicitly disclaim identifying the psychological process.
Elicited priors and researcher-fixed conditionals, interpreting certainty
responses, task wording and the sparse individual record remain material
assumptions. They discuss both an omitted dependence and conservative
human revision as possible explanations of mismatch. These observations
support a bounded judgment-model comparison, not a normative theorem
or a population claim that people always use Bayes nets.
\cite[accepted manuscript, pp.~31--36]{HarrisHahnMadsenHsu2016Expert}

\subsection{Consistency, credibility and guilt judgments in a trial task}
\label{sec:p16-rat42-legal-study}

Connor Desai, Reimers and Lagnado used an assault-trial transcript to
study whether inconsistency changes witness credibility and guilt judgments.
Study~1 reports 126 US/UK online recruits, 64 female, mean age
\(29.12\) with SD \(10.80\), paid \(\pounds2.40\), randomly split
among six conditions. Testimony consistency had two levels and the
defendant's prior-conviction information three. Participants read 20
evidential statements in a common order and repeatedly rated guilt and
victim/defendant credibility, then gave final ratings. Repeated statements
are not independent participants. Average completion time was 31 minutes.
\cite[pp.~626--628]{ConnorDesaiReimersLagnado2016Credibility}

Consistent testimony produced mean final guilt rating \(72.20\)
(SD \(26.49\)), versus \(58.98\) (SD \(29.03\)) for inconsistent
testimony, on the 0--100 scale; the reported ANOVA was
\(F(1,98)=12.35,p<.001\). Victim credibility likewise differed,
\(F(1,98)=16.97,p<.001\); defendant-credibility and prior-conviction
comparisons were not significant. The report gives 126 participants
but residual degrees of freedom 98 without explaining the corresponding
analysis-count discrepancy. The exact analyzed denominator for this
ANOVA is therefore not silently supplied.

Study~2 recruited 137 paid online participants, 65 female, mean age
\(32.52\), SD \(10.67\), randomly split between consistency conditions.
It elicited each participant's conditional probabilities \emph{after}
their posterior judgments and used them to parameterize a network of
selected testimony, intoxication, credibility and provocation. A
\(.50\) intoxication prior was fixed for all model calculations.
The model's provocation posterior served as a proxy for guilt, not as
an identical guilt variable. Its correlation with observed guilt ratings
was \(r=-.37,p<.001\), with reported \(R^2=.14\).
The observed guilt means \(73.87\) versus \(68.34\) differed in the
same direction as Study~1 but not significantly,
\(t(135)=-1.41,p>.05\). Both victim- and defendant-credibility
contrasts were significant. No numerical confidence interval for these
correlations is given in the selected report.
\cite[pp.~629--631]{ConnorDesaiReimersLagnado2016Credibility}

The authors note that the model includes only some testimony, omits
important defendant evidence, and uses provocation as a proxy. Parameter
elicitation after the outcomes is not held-out prediction. The trial
material is richer than a bare hypothesis-report problem but does not
make this online sample representative of actual juries, nor identify
a unique mental representation. The nonsignificant Study~2 guilt contrast
and the unaccounted Study~1 analysis denominator limit the stronger
claim that all effects were replicated. The paper supports a specific
modeling and judgment pattern, not legal advice or a verdict rule.

\section{Related treatments, successful use and failure cases}
\label{sec:p16-rat42-comparisons-and-use}

The probability operations agree with
\cref{sec:p15-finite-probability,sec:p15-bayesian-updating-prediction};
this source adds an analysis of choosing variables, modeling testimony,
and evaluating argumentation. The dependent-source example explains
why dependence is not synonymous with duplicated data. The existing
P16 structural-causal chapters add equations and intervention invariance:
\(P(Y\mid\operatorname{do}(X=x))\) is not licensed merely by a
Bayes-net arrow or by conditioning. The exact distinction and worked
counterexample are
\cref{def:p16-causal-intervention,thm:p16-causal-confounded-example}.

P19's empirical framework asks which task, population, information,
response measure and model comparison justify a cognitive interpretation.
Its suppression treatment concerns interpreted premises, whereas the
material-premise theorem above holds with fixed meanings. A description
of a reasoning error and a formal invalidity claim need not be the same
claim. Conversely reinterpreting every disagreement post hoc would make
a cognitive account untestable; representation choices need independent
constraints or predictive tests.

\begin{hbexample}{def:p16-rat42-success-failure}{Task and property lookup cases}
For calculating the effect of two witness reports, specify whether the
reliability variables are separate or shared and whether reports are
conditionally independent given their parents. The reliability theorem
then returns different justified posteriors under equal one-report
marginals. Using a single ``independent evidence is always better''
rule fails in its displayed low-\(a\), low-\(r\) case.

For evaluating an unsuccessful search argument, give the count prior,
the truth/count relation, the search-success relation and the screening
assumption. The covariance identity distinguishes confirmation,
disconfirmation and no information. Counting a search failure as
disconfirmation without a sensitive search model fails when all
\(f_i\)'s are equal.

For a cognitive claim, compare the participant-level elicited inputs and
judgments under the actual task. A close condition-mean fit can be a
success for a bounded computational-level model while failing to establish
individual accuracy, a mental algorithm or population generalization.
For a causal query, use a causal model's intervention semantics; merely
reversing a Markov-equivalent graph does not resolve that query.
\end{hbexample}

The source's outlook points to decision-making applications, ranking
theory and imprecise-probability independence. Those are extensions and
comparisons, not a claim that ordinary probability tables, ranking values
and sets of probabilities share one numerical semantics. Its references
to Spohn, Cozman and Halpern retain that cited boundary; it proposes
constructing and comparing analogous models rather than proving every
other framework here. The notes also locate the semigraphoid and
d-separation background, explain the variable/value convention and root
factor, and mention attempts at automated variable discovery. None
removes the need to justify the model of the actual reasoning task.
\cite[Sec.~5 and notes~1--10, pp.~262--264]{Hartmann2021BayesNets}
